\chapter{التأثيرات الإلكترونية والوسائط التفاعلية}\label{ch:b1:electronic-effects}

حمض ثلاثي فلورو الإيثانويك، \ce{CF3COOH}، أقوى حمضيةً من حمض الإيثانويك،
\ce{CH3COOH}، بنحو عشرة آلاف مرة، مع أن الرابطة O--H الحمضية
تبعد ثلاث روابط عن ذرات الفلور. والفينول أكثر حمضيةً من الإيثانول بمليون
مرة، وجزيء الأمونيا المرتبط بحلقة بنزين ---
الأنيلين --- أضعف قاعديةً بمليون مرة من جزيء أمونيا
مرتبط بمجموعة ميثيل. والكيمياء العضوية مليئة بمثل هذه التباينات،
ويفسّر معظمَها تأثيران: جذب الذرات الكهرسلبية على امتداد
الروابط $\sigma$، وانتشار الأزواج الإلكترونية على الأنظمة $\pi$.
والتأثيران نفساهما يقرّران أي الروابط تنكسر، وأي
الوسائط تتكوّن وأين تهاجم الكواشف: إنهما قواعد لغة
آليات التفاعل في الفصول التالية.

\begin{recall}
\omterm{def:b1:periodicity:electronegativity}{الكهرسلبية} وتغيّراتها (\cref{ch:b1:periodicity})؛ \omterm{def:b1:lewis-resonance-vsepr:lewis}{وبنى
لويس}، \omterm{def:b1:lewis-resonance-vsepr:formal-charge}{والشحنات الشكلية} والرنين (\cref{ch:b1:lewis-resonance-vsepr})؛
و$\mathrm{p}K_a$ وقوة الأحماض والقواعد
(\cref{ch:b1:predominant-reaction}). ورسم المجلد المدرسي (الصف الثاني عشر)
أسهمًا منحنية من \omterm{def:b1:electronic-effects:nucleophile}{محب للنواة} إلى \omterm{def:b1:electronic-effects:nucleophile}{محب للإلكترونات}.
\end{recall}

\section{التأثيرات التحريضية}

\begin{definition}[التأثير التحريضي]\label{def:b1:electronic-effects:inductive}
\emph{التأثير التحريضي}\index{تأثير تحريضي} هو انزياح
الكثافة الإلكترونية على امتداد الروابط $\sigma$ نحو ذرة أو
مجموعة كهرسلبية. وللمجموعة الجاذبة (F، Cl، O، \ce{NO2}، \ce{CF3}) تأثير $-I$؛
أما مجموعات الألكيل والذرات المشحونة سلبًا فتدفع الإلكترونات، وهذا تأثير $+I$.
ويضعف التأثير بسرعة مع عدد الروابط الفاصلة.
\end{definition}

\begin{proposition}[التأثير التحريضي والحمضية]\label{prop:b1:electronic-effects:inductive-pka}
المجموعة الجاذبة القريبة من الموقع الحمضي لحمض كربوكسيلي
تثبّت أيون الكربوكسيلات بنشر شحنته السالبة،
وتخفض $\mathrm{p}K_a$: فكلما كثرت المجموعات الجاذبة واقتربت، كان
الحمض أقوى.
\end{proposition}

\begin{proof}
تقيس $\mathrm{p}K_a$ موضع التوازن $\mathrm{RCOOH} \rightleftharpoons
\mathrm{RCOO^-} + \ce{H+}$. والمجموعة التي تسحب الكثافة الإلكترونية بعيدًا عن
الكربوكسيلات تنشر شحنتها على ذرات أكثر، فتخفض طاقته
أكثر مما تخفض طاقة الحمض المتعادل؛ فينزاح التوازن نحو
اليمين وتزداد $K_a$. وتؤكد ذلك السلسلةُ المقيسة.
\end{proof}

\begin{omfigure}
\begin{tikzpicture}
\begin{axis}[
  xbar, width=0.78\linewidth, height=5.2cm, xmin=0, xmax=5.4,
  symbolic y coords={TFA,TCA,DCA,MCA,AA}, ytick=data,
  yticklabels={\ce{CF3COOH}, \ce{CCl3COOH}, \ce{CHCl2COOH}, \ce{CH2ClCOOH}, \ce{CH3COOH}},
  xlabel={\foreignlanguage{arabic}{$\mathrm{p}K_a$ عند \qty{25}{\celsius}}}, bar width=9pt,
  nodes near coords, nodes near coords align={horizontal},
  every node near coord/.append style={font=\scriptsize},
  tick label style={font=\small}, label style={font=\small}, enlarge y limits=0.15]
\addplot[fill=omDef!45, draw=omDef] coordinates {(0.52,TFA) (0.51,TCA) (1.35,DCA) (2.87,MCA) (4.76,AA)};
\end{axis}
\end{tikzpicture}

{\small أحماض الإيثانويك المهلجنة. كل ذرة كلور تخفض
$\mathrm{p}K_a$؛ ومع ثلاث ذرات يصير الحمض أقوى بنحو عشرة آلاف مرة
من حمض الإيثانويك. وحمضا ثلاثي فلورو الإيثانويك وثلاثي كلورو الإيثانويك
متساويان في القوة في حدود عدم يقين القياسات على مثل هذه \omterm{def:b1:predominant-reaction:strong-acid}{الأحماض
القوية}.}
% ledger: pka:ethanoic, pka:chloroethanoic, pka:dichloroethanoic, pka:trichloroethanoic, pka:trifluoroethanoic
\end{omfigure}

وتعمل مجموعات الألكيل في الاتجاه المعاكس، قليلًا: فحمض البروبانويك ($\mathrm{p}K_a
= 4.89$) أضعف بقليل من حمض الإيثانويك (4.76)، وحمض الميثانويك
(3.74)، الخالي من مجموعة ألكيل، أقوى من كليهما.
% ledger: pka:propanoic, pka:methanoic

\section{التأثيرات الميزوميرية}

حين ترتبط ذرة تحمل \omterm{def:b1:lewis-resonance-vsepr:lewis}{زوجًا حرًا}، أو شحنة، بنظام $\pi$،
يمكن أن تتشارك إلكتروناتها مع ذلك النظام؛ وتصف \omterm{def:b1:lewis-resonance-vsepr:resonance}{بنى الرنين}
(\cref{ch:b1:lewis-resonance-vsepr}) هذه المشاركة.

\begin{definition}[التأثير الميزوميري، المجموعات المانحة والساحبة]\label{def:b1:electronic-effects:mesomeric}
\emph{التأثير الميزوميري}\index{تأثير ميزوميري} هو انزياح
الكثافة الإلكترونية عبر نظام $\pi$ مترافق، وتبيّنه \omterm{def:b1:lewis-resonance-vsepr:resonance}{بنى
الرنين}. و\emph{المجموعة المانحة}\index{مجموعة مانحة} تعطي \omterm{def:b1:lewis-resonance-vsepr:lewis}{زوجًا حرًا}
للنظام، وهذا تأثير $+M$ (\ce{-OH}، \ce{-O^-}، \ce{-NH2}، \ce{-OR})؛
و\emph{المجموعة الساحبة}\index{مجموعة ساحبة} تسحب منه إلكترونات عبر
رابطة $\pi$ مع ذرة كهرسلبية، وهذا تأثير $-M$ (\ce{-NO2}،
\ce{-C=O}، \ce{-C#N}).
\end{definition}

\begin{omfigure}
\setchemfig{atom sep=2em}
\schemestart
\chemfig{CH_3-C(=[1]O)-[7]O^{-}}
\arrow{<->}
\chemfig{CH_3-C(-[1]O^{-})=[7]O}
\schemestop
\hspace{2.5em}
\schemestart
\chemfig{*6(=-=(-O^{-})-=-)}
\arrow{<->}
\chemfig{*6((=O)-=-\chemabove{}{\scriptstyle\ominus}-=-)}
\schemestop

\medskip
\schemestart
\chemfig{H_3C-C(=[1]O)-[7]NH_2}
\arrow{<->}
\chemfig{H_3C-C(-[1]O^{-})=[7]NH_2^{+}}
\schemestop

{\small \omterm{def:b1:lewis-resonance-vsepr:resonance}{بنى الرنين} لأيون الإيثانوات (الشحنة تتقاسمها
ذرتا أكسجين متكافئتان)، ولأيون الفينوكسيد (الشحنة منتشرة في
الحلقة؛ ومُثّلت واحدة من بنى الحلقة الثلاث)، ولأميد (\omterm{def:b1:lewis-resonance-vsepr:lewis}{الزوج
الحر} للنيتروجين مشترك مع C=O، وهذا ما يجعل الأميدات \omterm{def:b1:predominant-reaction:strong-acid}{قواعد ضعيفة}).}
\end{omfigure}

\begin{proposition}[التأثير الميزوميري والحمضية]\label{prop:b1:electronic-effects:mesomeric-pka}
الحمض الذي تنشر قاعدته المرافقة شحنتها بالرنين أقوى
من الحمض الذي لا تستطيع قاعدته ذلك: فالأحماض الكربوكسيلية ($\mathrm{p}K_a$ نحو 4--5)
أقوى بكثير من الكحولات (نحو 16)، ويقع الفينول (10.0) بينهما.
\omterm{def:b1:electronic-effects:mesomeric}{والمجموعة الساحبة} في الموضع \textit{ortho} أو \textit{para} بالنسبة إلى OH
في فينول تقوّيه أكثر مما تقوّيه في الموضع \textit{meta}.
\end{proposition}

\begin{proof}
في الألكوكسيد تستقر الشحنة على ذرة أكسجين واحدة؛ وفي الكربوكسيلات
تتقاسمها ذرتا أكسجين؛ وفي الفينوكسيد تتقاسمها ذرة الأكسجين وثلاث
ذرات كربون من الحلقة، أقل كهرسلبيةً من الأكسجين، فالكسب أصغر.
ومجموعة \ce{NO2} على الكربون \textit{ortho} أو \textit{para} تستطيع أن تأخذ
الشحنة بنفسها عبر \omterm{def:b1:lewis-resonance-vsepr:resonance}{بنية رنين} فيها \ce{C=N+(-O^-)-O^-}؛ أما في
\textit{meta}، فلا توجد \omterm{def:b1:lewis-resonance-vsepr:resonance}{بنية رنين} تنقل الشحنة إلى الكربون
الذي يحمل \ce{NO2}، ولا يبقى إلا \omterm{def:b1:electronic-effects:inductive}{التأثير التحريضي}. والقيم
المقيسة: 2-نيتروفينول 7.23، و4-نيتروفينول 7.16، و3-نيتروفينول 8.36،
مقابل 9.99 للفينول و15.9 للإيثانول.
\end{proof}
% ledger: pka:ethanol, pka:phenol, pka:2-nitrophenol, pka:3-nitrophenol, pka:4-nitrophenol

والاستدلال نفسه ينطبق على القواعد. فالميثيل أمين ($\mathrm{p}K_a$ لأيون
الأمونيوم الموافق 10.66) قاعدة أقوى من الأمونيا (9.25، \cref{ch:b1:predominant-reaction})،
لأن مجموعة الميثيل تدفع الإلكترونات نحو النيتروجين؛ وفي الأنيلين يكون \omterm{def:b1:lewis-resonance-vsepr:lewis}{الزوج الحر}
مشتركًا مع الحلقة ويكون أيون الأمونيوم أكثر حمضيةً بكثير (4.60).
ويحتفظ البيريدين (5.23) بزوجه الحر، لكن على نيتروجين مثبَّت في حلقة
مستوية بطابع \textit{s} أكبر، فيكون أقل إتاحةً منه في الأمين.
% ledger: pka:methylammonium, pka:anilinium, pka:pyridinium, dfg:NH4+_ao

\section{كسر رابطة: الوسائط التفاعلية}

\begin{definition}[الانشطار المتجانس والانشطار غير المتجانس]\label{def:b1:electronic-effects:cleavage}
تنكسر \omterm{def:b1:lewis-resonance-vsepr:lewis}{الرابطة التساهمية} \emph{بانشطار متجانس}\index{انشطار متجانس} حين تحتفظ كل ذرة
بواحد من إلكتروني الرابطة، فينتج جذران حران،
و\emph{بانشطار غير متجانس}\index{انشطار غير متجانس} حين تحتفظ ذرة واحدة بالإلكترونين كليهما، فينتج
كاتيون وأنيون.
\end{definition}

\begin{definition}[الوسائط التفاعلية]\label{def:b1:electronic-effects:intermediates}
\emph{الوسيط التفاعلي}\index{وسيط تفاعلي} نوع كيميائي
يتكوّن في خطوة من آلية ويُستهلك في خطوة لاحقة، وهو أشد
تفاعليةً من أن يتراكم (\cref{ch:b1:elementary-steps}).
و\emph{الكاتيون الكربوني}\index{كاتيون كربوني} فيه كربون بثلاث روابط
وشحنة موجبة (ستة \omterm{def:b1:quantum-numbers:configuration}{إلكترونات تكافؤ})؛
و\emph{الأنيون الكربوني}\index{أنيون كربوني} فيه كربون بثلاث روابط \omterm{def:b1:lewis-resonance-vsepr:lewis}{وزوج حر}
وشحنة سالبة؛ و\emph{الجذر الحر}\index{جذر حر} ذرة فيها
إلكترون غير مزدوج.
\end{definition}

\begin{definition}[صنف الكربون]\label{def:b1:electronic-effects:carbon-class}
الكربون المرتبط بكربون واحد آخر \emph{كربون
أولي}\index{كربون أولي}، والمرتبط باثنين \emph{كربون
ثانوي}\index{كربون ثانوي}، والمرتبط بثلاثة \emph{كربون
ثالثي}\index{كربون ثالثي}. ويأخذ \omterm{def:b1:electronic-effects:intermediates}{الكاتيون الكربوني} أو \omterm{def:b1:electronic-effects:intermediates}{الجذر الحر} أو
الهالوجينو ألكان صنفَ الكربون المعني.
\end{definition}

\begin{omfigure}
\begin{tikzpicture}[font=\small]
  % carbocation: planar, empty p orbital
  \begin{scope}
    \fill[omProp!15, draw=omProp] (0,0.15) ellipse (0.18 and 0.75);
    \fill[omProp!15, draw=omProp] (0,-0.15) ellipse (0.18 and 0.75);
    \draw[thick] (0,0) -- (-1.0,-0.25) node[left] {R};
    \draw[thick] (0,0) -- (0.95,-0.35) node[right] {R};
    \draw[thick] (0,0) -- (0.35,0.45) node[above right] {R};
    \node[fill=white, inner sep=1pt] at (0,0) {\ce{C+}};
    \node at (0,-1.75) {\foreignlanguage{arabic}{كاتيون كربوني: مستوٍ،}};
    \node at (0,-2.15) {\foreignlanguage{arabic}{مدار \babelsublr{\textit{p}} فارغ}};
  \end{scope}
  % radical: nearly planar, one electron
  \begin{scope}[xshift=4.6cm]
    \draw[omMeth] (0,0.15) ellipse (0.18 and 0.75);
    \draw[omMeth] (0,-0.15) ellipse (0.18 and 0.75);
    \fill (0,0.55) circle (0.05);
    \draw[thick] (0,0) -- (-1.0,-0.25) node[left] {R};
    \draw[thick] (0,0) -- (0.95,-0.35) node[right] {R};
    \draw[thick] (0,0) -- (0.35,0.45) node[above right] {R};
    \node[fill=white, inner sep=1pt] at (0,0) {C};
    \node at (0,-1.75) {\foreignlanguage{arabic}{جذر حر: شبه مستوٍ،}};
    \node at (0,-2.15) {إلكترون واحد};
  \end{scope}
  % carbanion: pyramidal with a lone pair
  \begin{scope}[xshift=9.2cm]
    \draw[thick] (0,0) -- (-0.95,-0.55) node[left] {R};
    \draw[thick] (0,0) -- (0.95,-0.55) node[right] {R};
    \draw[thick] (0,0) -- (0.2,-0.9) node[below] {R};
    \fill[omDef!15, draw=omDef] (0,0.55) ellipse (0.22 and 0.45);
    \fill (-0.07,0.65) circle (0.045); \fill (0.07,0.65) circle (0.045);
    \node[fill=white, inner sep=1pt] at (0,0) {\ce{C-}};
    \node at (0,-1.75) {\foreignlanguage{arabic}{أنيون كربوني: هرمي،}};
    \node at (0,-2.15) {زوج حر};
  \end{scope}
\end{tikzpicture}

{\small أشكال \omterm{def:b1:electronic-effects:intermediates}{الوسائط التفاعلية} الثلاثة للكربون (R: H أو
مجموعة ألكيل). يمكن مهاجمة \omterm{def:b1:electronic-effects:intermediates}{الكاتيون الكربوني} المستوي من أيٍّ من وجهيه، وهي
نقطة مهمة للكيمياء الفراغية في
\cref{ch:b1:nucleophilic-substitution}.}
\end{omfigure}

\begin{proposition}[ثبات الكاتيونات الكربونية]\label{prop:b1:electronic-effects:carbocation-order}
تثبت \omterm{def:b1:electronic-effects:intermediates}{الكاتيونات الكربونية} بالمجموعات التي تعطيها إلكترونات: \babelsublr{\foreignlanguage{arabic}{ثالثي}
$>$ \foreignlanguage{arabic}{ثانوي} $>$ \foreignlanguage{arabic}{أولي} $>$ \foreignlanguage{arabic}{ميثيلي}}؛ \omterm{def:b1:electronic-effects:intermediates}{والكاتيون الكربوني} المجاور لرابطة ثنائية
(أليلي) أو لحلقة بنزين (بنزيلي) يزداد ثباتًا بالرنين،
والمجاور لذرة تحمل \omterm{def:b1:lewis-resonance-vsepr:lewis}{زوجًا حرًا} (\ce{R-O-CH2+}) يزداد ثباتًا
أكثر.
\end{proposition}

\begin{proof}
كل مجموعة ألكيل تدفع الكثافة الإلكترونية نحو الكربون الفقير
بالإلكترونات (تأثير $+I$، وتآثر ذو صلة لروابطها C--H مع المدار
الفارغ، يصفه مجلد السنة الثانية). وللكاتيون الأليلي
\omterm{def:b1:lewis-resonance-vsepr:resonance}{بنيتا رنين} تضعان الشحنة على أحد الكربونين الطرفيين؛ وللكاتيون
البنزيلي أربع، ثلاث منها في الحلقة؛ ويكوّن \omterm{def:b1:lewis-resonance-vsepr:lewis}{زوج حر} للأكسجين
رابطة رابعة مع الكربون الكاتيوني، \ce{R-O+=CH2}، وفيها لكل ذرة
ثمانية.
\end{proof}

وتتبع \omterm{def:b1:electronic-effects:intermediates}{الأنيونات الكربونية} الترتيب المعاكس (\babelsublr{\foreignlanguage{arabic}{ميثيلي} $>$ \foreignlanguage{arabic}{أولي} $>$ \foreignlanguage{arabic}{ثانوي}
$>$ \foreignlanguage{arabic}{ثالثي}}) وتثبّتها المجموعات الساحبة؛ أما \omterm{def:b1:electronic-effects:intermediates}{الجذور الحرة} فتتبع
ترتيب \omterm{def:b1:electronic-effects:intermediates}{الكاتيونات الكربونية} نفسه، بوضوح أقل.

\section{الأسهم المنحنية، والمحبات للنواة والمحبات للإلكترونات}

\begin{definition}[الأسهم المنحنية]\label{def:b1:electronic-effects:curly-arrow}
في \omterm{def:b1:elementary-steps:elementary-step}{آلية تفاعل}، يبيّن \emph{السهم المنحني}\index{سهم منحن}
حركة زوج إلكتروني: يبدأ من \omterm{def:b1:lewis-resonance-vsepr:lewis}{زوج حر} أو من رابطة وينتهي
على ذرة (حيث يتكوّن \omterm{def:b1:lewis-resonance-vsepr:lewis}{زوج حر} أو رابطة جديدة) أو بين ذرتين (رابطة
جديدة). و\emph{السهم ذو نصف الرأس}\index{سهم ذو نصف رأس} يبيّن
حركة إلكترون واحد، في التفاعلات الجذرية.
\end{definition}

\begin{method}[رسم الأسهم المنحنية]\label{met:b1:electronic-effects:curly-arrows}
\begin{enumerate}
  \item ابدأ من الإلكترونات: \omterm{def:b1:lewis-resonance-vsepr:lewis}{زوج حر}، أو رابطة $\pi$ أو رابطة $\sigma$،
        ولا تبدأ أبدًا من شحنة موجبة أو من ذرة بلا إلكترونات.
  \item انتهِ حيث يذهب الزوج: رابطة جديدة بين الذرتين، أو
        \omterm{def:b1:lewis-resonance-vsepr:lewis}{زوج حر} على الذرة الأكثر كهرسلبيةً حين تنكسر رابطة.
  \item عُدّ الشحنات بعد كل خطوة: الذرة التي تعطي \omterm{def:b1:lewis-resonance-vsepr:lewis}{زوجًا حرًا} لرابطة
        جديدة تكسب $+1$، والذرة التي تستقبل زوجًا بوصفه \omterm{def:b1:lewis-resonance-vsepr:lewis}{زوجًا حرًا}
        تكسب $-1$؛ والشحنة الكلية محفوظة.
  \item لا تتجاوز أبدًا الثمانية على C وN وO: حين يصل زوج، يجب أن
        يغادر زوج آخر.
\end{enumerate}
\end{method}

\begin{omfigure}
\setchemfig{atom sep=2.2em}
\schemestart
\chemfig{H_3C-C(-[2]H)(-[6]H)-[@{b}]@{br}Br}
\arrow{->[انشطار غير متجانس]}[,1.5]
\chemfig{H_3C-C(-[2]H)(-[6]H)^{+}}\+\chemfig{Br^{-}}
\schemestop
\chemmove{\draw[->, shorten <=2pt, shorten >=2pt](b) .. controls +(90:6mm) and +(90:6mm) .. (br);}

\medskip
\schemestart
\chemfig{H@{o}O^{-}}\+\chemfig{@{h}H-@{oc}O-H}
\arrow{<=>}
\chemfig{H_2O}\+\chemfig{HO^{-}}
\schemestop
\chemmove{\draw[->, thick, shorten <=2pt, shorten >=2pt](o) .. controls +(-80:8mm) and +(-100:8mm) .. (h);}
\par\vspace{5mm}

{\small \omterm{def:b1:electronic-effects:curly-arrow}{الأسهم المنحنية}. في الأعلى: \omterm{def:b1:electronic-effects:cleavage}{انشطار غير متجانس} للرابطة C--Br، يذهب الزوج
إلى البروم. في الأسفل: انتقال بروتون من الماء إلى الهيدروكسيد، مكتوبًا على أنه
هجوم \omterm{def:b1:lewis-resonance-vsepr:lewis}{زوج حر} للقاعدة على الهيدروجين (تنكسر عندئذ الرابطة O--H في
الماء ويبقى زوجها على الأكسجين؛ والسهم الثاني متروك لحلّه في
\cref{exo:b1:electronic-effects:4}).}
\end{omfigure}

\begin{definition}[المحب للنواة، المحب للإلكترونات، المجموعة المغادرة]\label{def:b1:electronic-effects:nucleophile}
\emph{المحب للنواة}\index{محب للنواة} نوع كيميائي يعطي
زوجًا إلكترونيًا لذرة غير الهيدروجين ليكوّن رابطة (\omterm{def:b1:lewis-resonance-vsepr:lewis-acid}{قاعدة
لويس})؛ و\emph{المحب للإلكترونات}\index{محب للإلكترونات} نوع كيميائي
يستقبل مثل هذا الزوج (\omterm{def:b1:lewis-resonance-vsepr:lewis-acid}{حمض لويس}، وكثيرًا ما يكون كربونًا يحمل شحنة
جزئية موجبة). و\emph{المحبة للنواة}\index{محبة للنواة} لنوع كيميائي
هي تفاعليته بوصفه محبًا للنواة، وتقاس \omterm{def:b1:rate-laws:order}{بثوابت السرعة}.
و\emph{المجموعة المغادرة}\index{مجموعة مغادرة} هي المجموعة التي ترحل
\omterm{def:b1:lewis-resonance-vsepr:lewis}{بالزوج الرابط} حين تنكسر رابطة مع الكربون انكسارًا غير متجانس.
\end{definition}

\begin{proposition}[المحبة للنواة والقاعدية]\label{prop:b1:electronic-effects:nucleophilicity-basicity}
من أجل \omterm{def:b1:electronic-effects:nucleophile}{المحبات للنواة} التي تهاجم بالذرة نفسها، تتبع \omterm{def:b1:electronic-effects:nucleophile}{المحبة للنواة}
القاعدية: \babelsublr{\ce{CH3O-} $>$ \ce{OH-} $>$ \ce{CH3COO-} $>$ \ce{H2O}}. ولا
يصح ذلك على امتداد عمود من الجدول الدوري في \omterm{def:b1:intermolecular-forces-solvents:solvent-classes}{المذيبات البروتونية}، حيث
تتفاعل \omterm{def:b1:electronic-effects:nucleophile}{المحبات للنواة} الكبيرة القابلة للاستقطاب الضعيفة الإذابة أسرع مع أنها
قواعد أضعف (\babelsublr{\ce{I-} $>$ \ce{Br-} $>$ \ce{Cl-} $>$ \ce{F-}})، ولا على
القواعد الضخمة، وهي \omterm{def:b1:electronic-effects:nucleophile}{محبات للنواة} ضعيفة.
\end{proposition}

\begin{proof}
تقيس القاعدية توازن الهجوم على بروتون،
وتقيس \omterm{def:b1:electronic-effects:nucleophile}{المحبة للنواة} سرعة الهجوم على كربون؛ وكلتاهما تزداد مع
إتاحة الزوج الإلكتروني، ومن هنا التوازي في حالة الذرة الواحدة.
وتتعلق السرعة أيضًا بالطاقة اللازمة لنزع المذيب عن
\omterm{def:b1:electronic-effects:nucleophile}{المحب للنواة} وبتشوّه سحابته الإلكترونية نحو
كربون مزدحم: فالأنيونات الصغيرة مثل \ce{F-} ترتبط بقوة بجزيئات \omterm{def:b1:intermolecular-forces-solvents:solvent-classes}{المذيب
البروتوني} (\cref{ch:b1:intermolecular-forces-solvents})،
والمجموعات الضخمة تعيق الاقتراب من الكربون لا من البروتون.
\end{proof}

\omterm{def:b1:electronic-effects:nucleophile}{المجموعة المغادرة} الجيدة \omterm{def:b1:predominant-reaction:strong-acid}{قاعدة ضعيفة}، ثابتة بعد أن تحمل الزوج:
\ce{I-}، \ce{Br-}، \ce{Cl-}، والماء، والسلفونات؛ أما \ce{OH-} و\ce{RO-}، وهما قاعدتان
قويتان، فمجموعتان مغادرتان رديئتان --- ولهذا يجب
تنشيط الكحولات قبل الاستبدال (\cref{ch:b1:alcohol-activation}).

\section{تمارين}

\begin{exercise}[$\star$]\label{exo:b1:electronic-effects:1}
صنّف كل كربون في 2-ميثيل البوتان وفي 2-برومو-2-ميثيل البروبان إلى
أولي أو ثانوي أو ثالثي أو رباعي (مرتبط بأربع ذرات كربون). وأي
\omterm{def:b1:electronic-effects:intermediates}{كاتيون كربوني} يتكوّن بأكبر سهولة من كل بروميد صيغته \ce{C4H9Br}؟
\end{exercise}

\begin{exercise}[$\star$]\label{exo:b1:electronic-effects:2}
احسب، انطلاقًا من قيم $\mathrm{p}K_a$، نسبة ثابتي
الحمضية لحمض كلورو الإيثانويك وحمض الإيثانويك، ولحمض ثلاثي فلورو الإيثانويك
وحمض الإيثانويك.
\end{exercise}

\begin{exercise}[$\star$]\label{exo:b1:electronic-effects:3}
ارسم \omterm{def:b1:lewis-resonance-vsepr:resonance}{بنى الرنين} لأيون الإيثانوات، ولأيون 4-نيتروفينوكسيد
(والشحنة على أكسجين من مجموعة النيترو) وللكاتيون
الأليلي \ce{CH2=CH-CH2+}.
\end{exercise}

\begin{exercise}[$\star$]\label{exo:b1:electronic-effects:4}
أكمل \omterm{def:b1:electronic-effects:curly-arrow}{الأسهم المنحنية} لانتقال البروتون في الشكل وللتفاعل
\ce{NH3 + H3O+ -> NH4+ + H2O}. وتحقق من الشحنات.
\end{exercise}

\begin{exercise}[$\star\star$]\label{exo:b1:electronic-effects:5}
رتّب تصاعديًا بحسب الحمضية: الإيثانول، والفينول، وحمض الإيثانويك،
و4-نيتروفينول، وحمض ثلاثي كلورو الإيثانويك، والماء ($\mathrm{p}K_a$
للمزدوجة $\ce{H2O}/\ce{OH-}$ يساوي 14.0). وعلّل كل خطوة بحجة تحريضية أو
ميزوميرية.
\end{exercise}

\begin{exercise}[$\star\star$]\label{exo:b1:electronic-effects:6}
رتّب تصاعديًا بحسب القاعدية: الأنيلين، والميثيل أمين، والأمونيا، والبيريدين.
وفسّر موضع الأنيلين \omterm{def:b1:lewis-resonance-vsepr:resonance}{ببنية رنين}.
\end{exercise}

\begin{exercise}[$\star\star$]\label{exo:b1:electronic-effects:7}
لماذا كان 3-نيتروفينول أضعف من 4-نيتروفينول لكنه يبقى أقوى من
الفينول؟ ارسم \omterm{def:b1:lewis-resonance-vsepr:resonance}{بنى الرنين} التي تعلّل الأمرين.
\end{exercise}

\begin{exercise}[$\star\star$]\label{exo:b1:electronic-effects:8}
رتّب \omterm{def:b1:electronic-effects:intermediates}{الكاتيونات الكربونية} \ce{CH3+}، و\ce{(CH3)3C+}، و\ce{CH2=CH-CH2+}،
و\ce{CH3CH2+}، و\ce{C6H5CH2+}، و\ce{CH3OCH2+}، وعلّل.
\end{exercise}

\begin{exercise}[$\star\star$]\label{exo:b1:electronic-effects:9}
عيّن \omterm{def:b1:electronic-effects:nucleophile}{المحب للنواة} \omterm{def:b1:electronic-effects:nucleophile}{والمحب للإلكترونات} في: \ce{CH3Br + OH-}؛
\ce{H2O + H+}؛ الإيثانال مع أيون السيانيد؛ \ce{BF3 + NH3}. وارسم
\omterm{def:b1:electronic-effects:curly-arrow}{الأسهم المنحنية}.
\end{exercise}

\begin{exercise}[$\star\star\star$]\label{exo:b1:electronic-effects:10}
يحتوي محلول على \qty{0.010}{mol/L} من حمض ثلاثي كلورو الإيثانويك. هل
تصح صيغة \omterm{def:b1:predominant-reaction:strong-acid}{الحمض الضعيف} $\mathrm{pH} = \frac12(\mathrm{p}K_a - \log C)$؟
احسب pH حسابًا صحيحًا، \omterm{def:b1:predominant-reaction:dissociation}{ودرجة التفكك}.
\end{exercise}

\begin{exercise}[$\star\star\star$]\label{exo:b1:electronic-effects:11}
تزداد قيم $\mathrm{p}K_a$ للكحولات المقيسة في الماء من الميثانول
(15.5) إلى الإيثانول (15.9) فالبروبان-2-أول (17.1) ف2-ميثيل البروبان-2-أول
(19.2). أعطِ تفسيرين، أحدهما \omterm{def:b1:electronic-effects:inductive}{بالتأثير التحريضي} لمجموعات
الألكيل على الألكوكسيد، والآخر بإذابة الألكوكسيد.
% ledger: pka:methanol, pka:ethanol, pka:2-propanol, pka:tBuOH
\end{exercise}

\begin{exercise}[$\star\star\star$]\label{exo:b1:electronic-effects:12}
فسّر لماذا لا تكون الأميدات قاعدية على النيتروجين ولا \omterm{def:b1:electronic-effects:nucleophile}{محبة للنواة}،
بينما الأمينات كذلك، مستعملًا \omterm{def:b1:lewis-resonance-vsepr:resonance}{بنية الرنين} للأميد.
وأي ذرة من الأميد تتبرتن في \omterm{def:b1:predominant-reaction:strong-acid}{الحمض القوي}؟
\end{exercise}

\section{مسألة: لماذا كان حمض ثلاثي فلورو الإيثانويك قويًا إلى هذا الحد}

\begin{problem}[{مسألة نهاية الأسبوع --- السلسلة التحريضية
للأحماض المهلجنة، والكربوكسيلات مقابل الألكوكسيدات، والنيتروفينولات
الثلاثة، وحساب بالتفاعل السائد، ونسبة
ثابتي الحمضية لحمض ثلاثي فلورو الإيثانويك وحمض الإيثانويك}]
\label{pb:b1:electronic-effects:1}
المعطيات ($\mathrm{p}K_a$ عند \qty{25}{\celsius}): حمض الإيثانويك 4.76،
كلورو الإيثانويك 2.87، ثنائي كلورو الإيثانويك 1.35، ثلاثي كلورو الإيثانويك 0.51،
ثلاثي فلورو الإيثانويك 0.52، الإيثانول 15.9، الفينول 9.99، 2-نيتروفينول 7.23،
3-نيتروفينول 8.36، 4-نيتروفينول 7.16؛ $\mathrm{p}K_e = 14.00$.

\medskip
\noindent\textbf{الجزء الأول --- السلسلة التحريضية.}
\begin{enumerate}
  \item احسب نسبة $K_a$ لكل خطوة من حمض الإيثانويك إلى
        حمض ثلاثي كلورو الإيثانويك.
  \item هل تأثير كل ذرة كلور هو نفسه؟ علّق.
  \item فسّر المنحى \omterm{def:b1:electronic-effects:inductive}{بالتأثير التحريضي}.
  \item لماذا لا يكاد يكون لذرة كلور على كربون سلسلة أطول، هو الأبعد
        عن COOH، أي تأثير؟
  \item الفلور أكثر كهرسلبيةً من الكلور. ماذا تُظهر
        مقارنة الحمضين الثلاثي الفلور والثلاثي الكلور، ولماذا
        يصعب قياس مثل قيم $\mathrm{p}K_a$ هذه بدقة؟
  \item في أي صورة يكون كل من هذه الأحماض عند pH يساوي 3؟
\end{enumerate}

\medskip
\noindent\textbf{الجزء الثاني --- الكربوكسيلات والألكوكسيدات.}
\begin{enumerate}[resume]
  \item ارسم \omterm{def:b1:lewis-resonance-vsepr:resonance}{بنيتي الرنين} لأيون الإيثانوات.
  \item بماذا تتنبآن بشأن طولي الرابطتين C--O في الأيون؟
  \item لماذا لا بنى متكافئة لأيون الإيثوكسيد؟
  \item احسب نسبة \omterm{def:b1:predominant-reaction:acidity-constant}{ثابت الحمضية} $K_a$ لحمض الإيثانويك إلى ثابت الإيثانول.
  \item ضع الفينول بين الاثنين وعلّل ذلك \omterm{def:b1:lewis-resonance-vsepr:resonance}{ببنى رنين}
        الفينوكسيد.
  \item أي من الإيثانول والفينول وحمض الإيثانويك يتفاعل تفاعلًا شبه تام
        مع محلول هيدروجين كربونات ($\mathrm{p}K_a$ يساوي 6.37 للمزدوجة
        $\ce{CO2,H2O}/\ce{HCO3-}$)؟
\end{enumerate}

\medskip
\noindent\textbf{الجزء الثالث --- النيتروفينولات.}
\begin{enumerate}[resume]
  \item رتّب النيتروفينولات الثلاثة والفينول بحسب الحمضية.
  \item ارسم \omterm{def:b1:lewis-resonance-vsepr:resonance}{بنية رنين} لأيون 4-نيتروفينوكسيد تكون فيها الشحنة على
        مجموعة النيترو.
  \item بيّن أنه لا توجد بنية كهذه لأيون 3-نيتروفينوكسيد.
  \item لماذا يبقى 3-نيتروفينول أكثر حمضيةً من الفينول؟
  \item 2-نيتروفينول أضعف قليلًا من 4-نيتروفينول مع أن كليهما
        يخضع \omterm{def:b1:electronic-effects:mesomeric}{للتأثير الميزوميري}؛ اقترح تفسيرًا يتعلق بمجموعة
        OH فيه وبمجموعة النيترو القريبة.
  \item عند pH يساوي 7.5، أي النيتروفينولات يكون في معظمه على صورة الأنيون؟ ولماذا
        يجعل ذلك 4-نيتروفينول كاشفًا لونيًا؟
\end{enumerate}

\medskip
\noindent\textbf{الجزء الرابع --- حمض ثلاثي فلورو الإيثانويك في الماء.}
\begin{enumerate}[resume]
  \item احسب pH لمحلول حمض ثلاثي فلورو الإيثانويك تركيزه \qty{0.010}{mol/L}
        (لا تفترض تفككًا ضعيفًا).
  \item احسب pH لمحلول حمض الإيثانويك تركيزه \qty{0.010}{mol/L}.
  \item اكتب تفاعل حمض ثلاثي فلورو الإيثانويك مع إيثانوات الصوديوم
        واحسب ثابته.
  \item اذكر النسبة $K_a(\ce{CF3COOH})/K_a(\ce{CH3COOH})$، على صورة قوة
        للعدد عشرة وعلى صورة عدد.
\end{enumerate}
\end{problem}
